\documentclass[11pt,a4paper]{article}
\usepackage{fontspec,amsmath,amssymb,amsthm,microtype}
\setmainfont{lmroman10-regular.otf}[BoldFont=lmroman10-bold.otf,ItalicFont=lmroman10-italic.otf,BoldItalicFont=lmroman10-bolditalic.otf]
\setmonofont{lmmono10-regular.otf}[Scale=MatchLowercase]
\usepackage[margin=25mm]{geometry}
\usepackage[unicode,hidelinks]{hyperref}
\setlength{\emergencystretch}{3em}
\setlength{\parskip}{0.4em}
\setlength{\parindent}{0pt}
\allowdisplaybreaks[2]
\raggedbottom
\title{\texorpdfstring{Bounded logarithmic tests and their entropy primitives}{Bounded logarithmic tests and their entropy primitives}}
\author{LangevinLean mathematical exposition}
\date{}
\begin{document}\maketitle
\noindent Source module: \texttt{Unadjusted\allowbreak{}Langevin/\allowbreak{}LSI/\allowbreak{}Bounded\allowbreak{}Log.\allowbreak{}lean}. The accompanying HTML pairs this mathematical exposition with the exact Lean statements and records the scope of the formalization.


The unbounded logarithm in entropy production is replaced by a smooth bounded function. The replacement has derivative between zero and one, converges to the identity, and admits an entropy primitive dominated by the actual KL integrand. These properties separately justify the spatial test, the time chain rule, and the limit in the entropy identity.


\section{\texorpdfstring{Choose the bounded test and its weight}{Choose the bounded test and its weight}}


For a natural number \(m\), write \(c=m+1>0\). Define

\[\chi_m(z)=c\arctan(z/c),\qquad w_m(z)=\frac1{1+(z/c)^2},\qquad\Phi_m(r)=\int_1^r\chi_m(\log u)du.\]

The positive offset in \(c\) prevents division by zero even for \(m=0\). The interval integral is oriented: when \(r<1\), its sign reverses. Lean defines these expressions for every real argument; the entropy claims below explicitly require \(r>0\).



{\raggedright\noindent Lean declarations: \texttt{bounded\allowbreak{}Log\allowbreak{}Slope}, \texttt{bounded\allowbreak{}Log\allowbreak{}Weight}, \texttt{bounded\allowbreak{}Entropy}.\par}

\section{\texorpdfstring{Differentiate and bound the test}{Differentiate and bound the test}}


The arctangent chain rule gives \(\chi_m(0)=0\) and \(\chi_m^{\prime}(z)=w_m(z)\). The denominator is at least one, so \(0\le w_m\le1\). The derivative bound implies that \(\chi_m\) is 1-Lipschitz, hence \(|\chi_m(z)|\le|z|\). Independently, \(|\arctan u|<\pi/2\) gives \(|\chi_m(z)|\le c\pi/2\). This last bound controls the time derivative integrand without any log-integrability requirement. Since arctangent has the sign of its argument, \(\chi_m\) does too.



\texttt{Cont\allowbreak{}Diff \ensuremath{\mathbb R} \ensuremath{\top}} records smoothness of the scalar test. \texttt{Has\allowbreak{}Deriv\allowbreak{}At} gives its exact derivative; the proof multiplies the derivative of arctangent by \(1/c\) and cancels the outer factor \(c\).



{\raggedright\noindent Lean declarations: \texttt{bounded\allowbreak{}Log\allowbreak{}Slope\_\allowbreak{}zero}, \texttt{bounded\allowbreak{}Log\allowbreak{}Slope\_\allowbreak{}cont\allowbreak{}Diff}, \texttt{bounded\allowbreak{}Log\allowbreak{}Slope\_\allowbreak{}has\allowbreak{}Deriv\allowbreak{}At}, \texttt{bounded\allowbreak{}Log\allowbreak{}Weight\_\allowbreak{}nonneg}, \texttt{bounded\allowbreak{}Log\allowbreak{}Weight\_\allowbreak{}le\_\allowbreak{}one}, \texttt{bounded\allowbreak{}Log\allowbreak{}Slope\_\allowbreak{}lipschitz}, \texttt{norm\_\allowbreak{}bounded\allowbreak{}Log\allowbreak{}Slope\_\allowbreak{}le}, \texttt{norm\_\allowbreak{}bounded\allowbreak{}Log\allowbreak{}Slope\_\allowbreak{}le\_\allowbreak{}const}, \texttt{bounded\allowbreak{}Log\allowbreak{}Slope\_\allowbreak{}nonneg}, \texttt{bounded\allowbreak{}Log\allowbreak{}Slope\_\allowbreak{}nonpos}.\par}

\section{\texorpdfstring{Recover the logarithm and its primitive}{Recover the logarithm and its primitive}}


As \(m\to\infty\), \(z/(m+1)\to0\). The derivative of arctangent at zero is one; its difference quotient therefore gives \(\chi_m(z)\to z\). At \(z=0\) the claim is exact at every index. Continuity of inversion at the positive denominator one gives \(w_m(z)\to1\).



The function \(u\mapsto\chi_m(\log u)\) is measurable and bounded, hence integrable on every finite interval, including intervals containing zero. When \(r>0\), continuity in a neighborhood of \(r\) and the fundamental theorem of calculus imply \(\Phi_m^{\prime}(r)=\chi_m(\log r)\).



{\raggedright\noindent Lean declarations: \texttt{bounded\allowbreak{}Log\allowbreak{}Slope\_\allowbreak{}tendsto}, \texttt{bounded\allowbreak{}Log\allowbreak{}Weight\_\allowbreak{}tendsto}, \texttt{bounded\allowbreak{}Log\_\allowbreak{}integrable}, \texttt{bounded\allowbreak{}Entropy\_\allowbreak{}has\allowbreak{}Deriv\allowbreak{}At}.\par}

\section{\texorpdfstring{Dominate the entropy endpoints}{Dominate the entropy endpoints}}


We have \(\Phi_m(1)=0\). For \(r\ge1\), both \(\log u\) and its test are nonnegative on \([1,r]\), giving \(\Phi_m(r)\ge0\). For \(0<r\le1\), the test is nonpositive on \([r,1]\), and reversing the interval again gives nonnegativity.



On the positive side, \(0\le\chi_m(\log u)\le\log u\). On the negative side, \(\log u\le\chi_m(\log u)\le0\); reversing the interval preserves the corresponding upper bound for the primitive. In both cases,

\[0\le\Phi_m(r)\le\int_1^r\log u\,du=r\log r-r+1.\]

Pointwise convergence of the test and domination by \(|\log u|\), integrable on the finite interval between 1 and \(r\), prove \(\Phi_m(r)\to r\log r-r+1\). These are exactly the hypotheses for dominated convergence against the finite-KL relative density later.



{\raggedright\noindent Lean declarations: \texttt{bounded\allowbreak{}Entropy\_\allowbreak{}one}, \texttt{bounded\allowbreak{}Entropy\_\allowbreak{}nonneg}, \texttt{bounded\allowbreak{}Entropy\_\allowbreak{}le\_\allowbreak{}kl\allowbreak{}Fun}, \texttt{bounded\allowbreak{}Entropy\_\allowbreak{}tendsto}.\par}

\section{\texorpdfstring{Identify the spatial gradient of the test}{Identify the spatial gradient of the test}}


For a differentiable scalar field \(f\), the Fréchet chain rule gives \(D(\chi_m\circ f)=w_m(f)Df\). Applying the Hilbert-space identification of linear derivatives with gradients yields

\[\nabla(\chi_m\circ f)(x)=w_m(f(x))\nabla f(x),\qquad\|\nabla(\chi_m\circ f)(x)\|\le\|\nabla f(x)\|.\]

The inequality uses \(0\le w_m\le1\). With \(f=\log r_t\), it supplies square-integrability of the bounded log test gradient from the actual Fisher bound.



{\raggedright\noindent Lean declarations: \texttt{gradient\_\allowbreak{}bounded\allowbreak{}Log\allowbreak{}Slope\_\allowbreak{}comp}, \texttt{norm\_\allowbreak{}gradient\_\allowbreak{}bounded\allowbreak{}Log\allowbreak{}Slope\_\allowbreak{}comp\_\allowbreak{}le}.\par}

\end{document}
